\subsection{Fazit}
Ziel dieses Feldkurses war das bestehende Netz 'Überwachung des Schwanderbärgli' zu modernisieren und zu vereinfachen.
Dabei sollten neue Ansätze präsentiert werden, welche jedoch dieselbe Genauigkeit und Vergleichbarkeit aufweisen.
In dieser Messkampagne wurden deutlich mehr GNSS-Messungen durchgeführt im Vergleich zu den vorjährigen Messungen.
Das erhoffte Ziel, das innere Netz 'Brunni' lediglich mittels GNSS-Messungen in
das Gesamtnetz zu integrieren konnte leider nicht erreicht werden.
Ein ausschlaggebender Grund war, dass der GDOP des Punktes 6011 nahezu permanent den Wert vier übersteigt.
In Anbetracht der Umliegenden Bewaldung und Abdeckung ist dies auch nicht sehr erstaunlich, wollte jedoch nicht unversucht bleiben.

\section{Reflexion}
\label{sec:reflexion}

Dies ist Platzhaltertext für das Reflexion. Hier werden die wichtigsten Erkenntnisse des
Berichts kurz zusammengefasst und ein Ausblick auf allfällige weitere Schritte gegeben.
